\documentclass[11pt]{amsart}

\usepackage[T1]{fontenc}
\usepackage[utf8]{inputenc}
\usepackage{amsmath,amssymb,amsthm}
\usepackage[margin=1in]{geometry}
\usepackage{booktabs}
\usepackage{enumitem}
\usepackage{hyperref}

\newcommand{\LL}{\Lambda}
\newcommand{\OO}{\mathbb{O}}
\newcommand{\ZZ}{\mathbb{Z}}
\newcommand{\FF}{\mathbb{F}}
\newcommand{\RR}{\mathbb{R}}
\newcommand{\Nrm}{N}

\newtheorem*{observation}{Observation}

\title[Mod-2 structure of $L/2L$: external proposal, verification, corrections]%
{Structural notes on $L/2L$ for\\
\emph{An order on the Leech lattice from a $\ZZ_3$-symmetric
triple-octonion product}\\
\large Subalgebra and ideal structure of the mod-2 quotient}

\author{Hermes\,$^{1}$ (proposed discoveries) and Claude Opus~4.7\,$^{2}$ (verification)}
\address{$^{1}$External contribution}
\address{$^{2}$Anthropic, acting at the direction of Jens K\"oplinger}

\date{24 May 2026}

\thanks{This supplement records two structural proposals about the
mod-2 quotient $L/2L$ communicated by an external contributor
identified as \emph{Hermes} (using OpenCode with DeepSeek), alongside an empirical
verification carried out in this project. Where the proposals were
correct or partially correct, the structure is recorded with
attribution; where corrections were required, the original claim is
reproduced verbatim and corrected forward. The supplement is companion material to the main paper
(\texttt{paper/main.tex}, v4); it is not part of the paper and is not
in its referee scope.}

\begin{document}
\maketitle

\begin{abstract}
The main paper's Appendix~A, Section~A.1 (``The mod-2-quotient reading
of Tables~A.2 and~A.3'') stops at the structural statement: under the
descended octonion product on $\bar L := L/2L$, the subspace
$V := \sigma(L\bar s)/2L$ is a $4$-dimensional left ideal and
$W := \sigma(Ls)/2L$ is a $4$-dimensional subalgebra, with
$\bar L = V \oplus W$.

External feedback from \emph{Hermes} (using OpenCode with DeepSeek) proposed two
finer structural descriptions: an explicit basis and multiplication
table for $W$ (with a~3-dim nilpotent ideal and a universal zero
divisor), and an explicit basis for~$V$ together with total isotropy
under the natural quadratic form on~$\bar L$. We reproduce both
proposals verbatim, then record the empirical verification carried out
in this project
(\texttt{python\_project/src/verify\_discovery1\_W\_subalgebra.py} and
\texttt{python\_project/src/verify\_discovery2\_V\_isotropic.py}).

Of the two proposals, one (Discovery~1, the $W$ multiplication table)
required two cell corrections and one structural re-interpretation;
the other (Discovery~2, $V$ totally isotropic) is correct in its
asserted form and is sharpened here in two new directions: $V$~is
in fact a \emph{two-sided} ideal of $\bar L$ (not only a left ideal,
as in the paper), and~$W$ is also totally isotropic under the natural
quadratic form (correcting the proposed ``mixed isotropy''
claim). Together, $V$ and $W$ form a Witt decomposition of $\bar L$
into a pair of maximal totally isotropic subspaces (Lagrangians) under
the plus-type quadratic form induced by the norm on~$L$.

The Lagrangian-polarization framing developed here is now recorded
at the end of Section~A.1 of the main paper.
\end{abstract}

% ===================================================================
\section{Setting}
% ===================================================================

Throughout, $L = D_8^+$ is the $E_8$ lattice in the
coordinate-symmetric placement, $\sigma = (1\;2)$ is the transposition
of two imaginary basis units of the paper, $L\bar s$ and~$Ls$ are
Wilson's sublattices, and $\bar L := L/2L \cong \FF_2^{\,8}$
is the mod-2 quotient with the descended (Z-bilinear $\to$ $\FF_2$-bilinear)
octonion product~$\bar\mu$. The paper's Appendix~A,
Section~A.1 establishes:

\begin{enumerate}[label=(\arabic*)]
\item $\bar L$ has an inherited $\FF_2$-bilinear product $\bar\mu$
  whose structure constants are the entries of Table~A.1 reduced
  modulo~$2$. The product is non-unital ($e_0 \notin L$) and
  non-commutative in the paper's $L$-basis $L_1,\ldots,L_8$.
\item $V := \sigma(L\bar s)/2L$ and $W := \sigma(Ls)/2L$ are
  $4$-dimensional $\FF_2$-subspaces of $\bar L$ with
  $\bar L = V \oplus W$.
\item Lemma~4.3 $\iff$ $V$ is a left ideal of $(\bar L, \bar\mu)$.
\item Lemma~4.4 $\iff$ $W$ is a subalgebra of $(\bar L, \bar\mu)$.
\end{enumerate}

The proposals below refine~(3) and~(4). All claims have been verified
by exact-arithmetic enumeration on the paper's Appendix-A basis;
verification scripts are in
\texttt{python\_project/src/verify\_discovery1\_W\_subalgebra.py}
and~\texttt{verify\_discovery2\_V\_isotropic.py}.

% ===================================================================
\section{Discovery 1 (Hermes): the multiplication table of $W$}
% ===================================================================

\subsection*{Original proposal (verbatim)}

\begin{quote}
\textbf{Discovery~1.}  $W = \sigma(Ls)/2L$ -- explicit subalgebra
identified. I computed the multiplication table of $W = \sigma(Ls)/2L$
as an $\FF_2$-algebra. This is the structural reason for Lemma~4.4.

Basis in $L/2L$ coordinates:
\begin{align*}
  w_0 &= (1,1,1,0,0,0,0,0) \text{ -- contains identity of } L/2L \\
  w_1 &= (0,1,0,1,1,0,1,0) \\
  w_2 &= (0,0,0,1,0,1,1,1) \text{ -- universal zero divisor} \\
  w_3 &= (0,0,0,0,1,0,1,1)
\end{align*}
Multiplication table:
\[
\begin{array}{c|cccc}
 & w_0 & w_1 & w_2 & w_3 \\ \hline
w_0 & w_0+w_1   & 0 & w_2 & 0 \\
w_1 & w_1+w_2+w_3 & 0 & 0   & w_2 \\
w_2 & 0         & 0 & 0   & 0 \\
w_3 & w_3       & w_2 & 0 & 0 \\
\end{array}
\]
Structural features:
\begin{itemize}
\item $N = \mathrm{span}\{w_1, w_2, w_3\}$ is a 3-dim nilpotent ideal,
  $N^2 = \mathrm{span}\{w_2\}$, $N^3 = 0$.
\item $w_2$ is a universal zero divisor.
\item The identity of $L/2L$ is in $W$ but NOT idempotent in $W$
  ($w_0^2 = w_0 + w_1$).
\item Only 6/64 basis products are non-zero.
\end{itemize}
This replaces the appendix tables with an abstract structural reason:
Lemma~4.4 holds because $W \in \mathrm{Subalg}(L/2L)$, provable
structurally rather than by enumerating 64 coefficients.
\end{quote}

\subsection*{Verification result}

The proposed basis of~$W$ is correct: the four vectors
$(w_0, w_1, w_2, w_3)$ are in $W$ and span exactly the $4$-dimensional
$\FF_2$-subspace $W = \sigma(Ls)/2L$.

The multiplication table requires two cell corrections, and one
structural claim is incorrect; the original framing as an
``$\FF_2$-algebra in the conventional sense'' overstates the
structure since the descended product on $\bar L$ is non-unital
(the octonion identity $e_0 \notin L$) and non-commutative in the
$L$-basis.

\paragraph{Actual multiplication table.}
Verified by direct computation of $\bar\mu(w_i, w_j)$ using the
structure constants in the paper's Appendix-A $L$-basis:
\[
\begin{array}{c|cccc}
\bar\mu(\cdot,\cdot) & w_0 & w_1 & w_2 & w_3 \\ \hline
w_0 & w_0 & 0   & w_2 & 0   \\
w_1 & w_1 & 0   & 0   & w_2 \\
w_2 & 0   & 0   & 0   & 0   \\
w_3 & w_3 & w_2 & 0   & 0   \\
\end{array}
\]
Six non-zero entries; every non-zero entry is a single basis vector.

\paragraph{Identity / idempotent claim.}
In $L$-basis coordinates, $w_0 = L_1 + L_2 + L_3$. Since
$L_2 + L_3 = (e_0 + e_1) + (e_0 - e_1) = 2 e_0 \in 2L$, this reduces
modulo~$2L$ to $w_0 \equiv L_1 = s$. The octonion identity~$e_0$ is not
in~$L$ (it has odd coordinate sum, so it lies outside~$D_8^+$), and
$\bar L$ has no multiplicative identity. The corrected statement: $w_0$
represents~$s$ modulo~$2L$ and is an \emph{idempotent in~$W$}
($w_0^2 = w_0$), not an inherited identity of~$\bar L$.

\paragraph{Structural features (corrected confirmations).}
With the actual table:
\begin{itemize}
\item $N := \mathrm{span}_{\FF_2}\{w_1, w_2, w_3\}$ is a $3$-dimensional
  two-sided ideal of~$W$.
\item $N^2 = \FF_2 w_2$ and $N^3 = 0$.
\item $w_2$ is a universal zero divisor of~$W$ (every product
  $\bar\mu(w_2, w) = 0$ for all $w \in W$, and $\bar\mu(w, w_2)$ is
  $0$ or $w_2$).
\item $W$ decomposes as $W = \FF_2 w_0 \oplus N$, with $w_0$
  idempotent.
\end{itemize}

\subsection*{Does this replace the appendix tables?}

Not by itself. The mod-2 quotient subsection (paper Section~A.1) already
converts Lemma~4.4 to ``$W$ is a subalgebra of $\bar L$''; the table
above describes the internal structure of~$W$ as an algebra, which
is a refinement. Verifying $W$~closes under $\bar\mu$ still requires
checking $4 \times 4 = 16$ basis products in~$\FF_2$~--- a reduction
from the appendix tables' $64$ integer products in $\ZZ$, but not an
elimination. A genuine structural argument would identify~$W$
intrinsically inside~$\bar L$ (kernel of a natural map, image of a
canonical construction, Peirce component, radical-like piece) without
referencing~$\sigma$ at all. That identification is not yet in hand.

% ===================================================================
\section{Discovery 2 (Hermes): $V$ totally isotropic left ideal}
% ===================================================================

\subsection*{Original proposal (verbatim)}

\begin{quote}
\textbf{Discovery~2.}  $V = \sigma(L\bar s)/2L$ is a totally isotropic
left ideal. $V$ basis:
\[ (0,1,1,0,0,0,0,0),\ (0,0,1,0,1,0,1,0),\ (0,0,0,1,0,0,1,1),\
   (0,0,0,0,1,1,0,1). \]
Properties verified:
\begin{itemize}
\item $V$ is a left ideal: $\bar\mu(L/2L, V) \subseteq V$
  (Lemma~4.3 structurally).
\item $V$ is totally isotropic: $q(v) = 0$ for all $v \in V$.
\item $V \oplus W = L/2L$.
\item $W$ has mixed isotropy ($8$ with $q=0$, $8$ with $q=1$).
\end{itemize}
\end{quote}

\subsection*{The quadratic form}

For an even unimodular lattice $L$ with norm form $\Nrm(v) =
\sum v_i^2$, the descended $\FF_2$-valued quadratic form
\[
  q \colon L/2L \to \FF_2, \qquad
  q(v + 2L) := \Nrm(v)/2 \pmod{2}
\]
is well-defined because $\Nrm(v + 2u) = \Nrm(v) + 4 \langle v, u\rangle
+ 4 \Nrm(u) \equiv \Nrm(v) \pmod{4}$, so the right-hand side is
$\sigma$-independent of the chosen representative. For $L = D_8^+$
this $q$ is the standard plus-type form on $\FF_2^{\,8}$: $136$
isotropic vectors (including~$0$) and $120$ anisotropic vectors out
of~$256$, Witt index~$4$.

\subsection*{Verification result}

The proposed basis of~$V$ is correct: the four vectors above are in
$V = \sigma(L\bar s)/2L$ and span exactly it.

The two correct claims and the four-by-four scoresheet:

\begin{center}
\begin{tabular}{lc}
\toprule
Claim & Verified \\
\midrule
$V$ is a left ideal of $\bar L$ & Yes (already paper Lemma~4.3) \\
$V$ is totally isotropic, $q \equiv 0$ on $V$ & Yes (all $16$ elements)\\
$V \oplus W = \bar L$ & Yes (already paper Section~A.1) \\
$W$ has mixed isotropy ($8/8$) & \textbf{No} -- $W$ is totally isotropic \\
\bottomrule
\end{tabular}
\end{center}

\paragraph{$W$ is totally isotropic (correction to proposed mixed
isotropy).}
Direct computation gives $q(w) = 0$ for all $16$ elements
$w \in W$. This is empirically clean: $V$ and $W$ are \emph{both}
$4$-dimensional totally isotropic subspaces of $\bar L$, and the
proposed $8/8$ split for~$W$ is incorrect under the natural
quadratic form $q(v) = \Nrm(v)/2 \pmod 2$. (The proposal may have used
a different form; under the standard one, $W$ is totally isotropic.)

\paragraph{Bonus structural fact: $V$ is also a right ideal.}
The paper records $V$ as a left ideal (the mod-2 form of Lemma~4.3).
Direct empirical test gives also
\[ \bar\mu(V, \bar L) \subseteq V, \]
so $V$ is in fact a \emph{two-sided ideal} of $\bar L$. This is a
genuine structural fact not appearing in the paper's mod-2 reading,
and the right-ideal direction does not follow from the left-ideal
direction (the algebra $\bar\mu$ is non-commutative in the
$L$-basis).

\subsection*{The Witt decomposition picture}

Putting the verified pieces together: under the plus-type quadratic
form $q$ on $\bar L \cong \FF_2^{\,8}$,
\begin{equation}\label{eq:witt}
  \bar L \;=\; V \oplus W,
\end{equation}
where:
\begin{itemize}
\item $V$ is a $4$-dimensional \emph{maximal totally isotropic
  two-sided ideal} of $(\bar L, \bar\mu)$;
\item $W$ is a $4$-dimensional \emph{maximal totally isotropic
  subalgebra} of $(\bar L, \bar\mu)$;
\item $V \cap W = \{0\}$, and $V, W$ together span~$\bar L$;
\item the Witt index of the plus-type form on $\FF_2^{\,8}$ is~$4$,
  so each of $V$ and $W$ is a maximal isotropic subspace
  (Lagrangian).
\end{itemize}

This is a clean orthogonal-geometry picture of what
$\sigma$ accomplishes structurally: a transversal decomposition of
$\bar L$ into a pair of complementary Lagrangians, one carrying the
two-sided-ideal role and the other carrying the subalgebra role.

\begin{observation}
$V$ and $W$ are not only orthogonally placed -- they are
geometrically related: under the bilinear form
$b(x, y) := q(x+y) - q(x) - q(y)$ associated to~$q$, the pairing
between $V$ and $W$ is non-degenerate (since their direct sum is
$\bar L$ and both are isotropic), so $V$ and $W$ are dual Lagrangians
of the orthogonal space~$(\bar L, q)$.
\end{observation}

% ===================================================================
\section{Discovery 1 vs Discovery 2: are the ideals the same?}
% ===================================================================

The user (and a reader) might reasonably ask whether the
$3$-dim nilpotent ideal $N \subset W$ from Discovery~1 is the same
object as the $4$-dim two-sided ideal $V \subset \bar L$ from
Discovery~2. \emph{They are not.} Their relationship is:

\begin{itemize}
\item $V$~lives in the complementary summand to~$W$ in the
  decomposition $\bar L = V \oplus W$;
\item $N$~lives \emph{inside}~$W$ (it is a sub-ideal of the
  subalgebra~$W$), and is therefore disjoint from~$V$ except at~$0$
  (empirical: $|N \cap V| = 1$).
\end{itemize}

The two ideals are structurally distinct components of the same
overall mod-2 picture: $V$ is the ``ideal half'' of the
$V \oplus W$ decomposition (it carries the left-multiplication
action of $\bar L$), and $N$ is the radical-like core of the
``subalgebra half''~$W$. Both are totally isotropic under~$q$;
$V$~is in addition a two-sided ideal of all of~$\bar L$, whereas
$N$~is only known to be a two-sided ideal of~$W$.

% ===================================================================
\section{Relevance to the paper's open question}
% ===================================================================

The paper's Section~7 records as open: \emph{a structural reason for
Lemma~\ref{lem:sigmaLs-closed}}, asking specifically for an
``identification of~$\sigma(Ls)$ with a left ideal, with the image
under a norm-preserving octonion automorphism, or with some
module-theoretic structure over~$L$, that would lift the mod-2
picture''. The structure documented here is the next layer past
Section~A.1: not the integer-level structural reason the open
question targets, but a sharper account of the mod-2 picture itself
--- specifically, the Witt-decomposition interpretation and the
internal multiplication table of~$W$. These do not close the open
question (the integer-level lift remains open) but they refine the
target: any future intrinsic identification of $\sigma(Ls)$ should
descend to recover the multiplication table of~$W$ recorded in
\S2, with $w_0$ idempotent and $w_2$ a universal zero divisor.

% ===================================================================
\section*{Attribution and provenance}
% ===================================================================

Both structural proposals (\textbf{Discovery~1} and \textbf{Discovery~2})
were communicated by an external contributor identified as
\emph{Hermes} (using OpenCode with DeepSeek). The proposals reached this project
through review feedback on the main paper's Appendix~A
mod-2-quotient subsection; the original wording is reproduced
verbatim in \S2 and~\S3.

The empirical verification (existence of~$V$ and~$W$ with the
proposed bases, the multiplication table of~$W$, total isotropy
of~$V$, total isotropy of~$W$ under the natural quadratic form,
the two-sided-ideal property of~$V$, and the structural separation
of~$N \subset W$ from~$V \subset \bar L$) was carried out in this
project by Claude Opus~4.7 (Anthropic) at the direction of Jens
K\"oplinger, using the scripts
\texttt{python\_project/src/verify\_discovery1\_W\_subalgebra.py}
and \texttt{verify\_discovery2\_V\_isotropic.py}.

Corrections to the original proposals (the two-cell error in
the $W$~multiplication table, the misidentification of $w_0$ with
``the identity of $L/2L$'', and the mistaken ``$W$ has mixed
isotropy'' claim) are recorded in \S2--\S3 with the verbatim
originals preserved.

% ===================================================================
\section*{Cross-references}
% ===================================================================

\begin{itemize}
\item Paper Appendix~A, Section~A.1 (``The mod-2-quotient reading of
  Tables~A.2 and~A.3''): the level at which the paper's structural
  account stops.
\item Paper~\S7, open question~(1): the integer-level structural
  reason for Lemma~4.4 (\emph{remains open}).
\item \texttt{python\_project/src/verify\_mod2\_quotient.py}: paper-side
  verification of $V$ as a left ideal and $W$ as a subalgebra.
\item \texttt{python\_project/src/verify\_discovery1\_W\_subalgebra.py}:
  verification of the $W$-multiplication table (with two-cell
  corrections to the proposal).
\item \texttt{python\_project/src/verify\_discovery2\_V\_isotropic.py}:
  verification of $V$ totally isotropic, two-sided ideal, $W$ also
  totally isotropic.
\item \texttt{paper/reviews/2026-05-24\_claude\_opus\_4\_7\_W\_concrete.md}:
  the project-internal record of the $W$ concrete example.
\item \texttt{prompt\_logs/155\_W\_subalgebra\_concrete.txt} and
  \texttt{156\_appendix\_renumber\_discovery\_record.txt}: prompt logs.
\end{itemize}

\end{document}
